\section{Significance of the Study}

This study holds significant value for multiple sectors, ranging from small-scale backyard farmers to major commercial piggery operators, as well as broader agricultural and technological stakeholders.

\noindent\textit{Swine Farmers and Caretakers}

By replacing manual feeding with a sensor-based dispensing and monitoring system, the proposed solution directly addresses feeding problems such as inconsistent rationing, excess feed distribution, feed wastage, and difficulty tracking actual feed intake per pen. The system measures the amount of feed dispensed and consumed through the load cell, allowing caretakers to compare current consumption with previous feeding records. When feed intake drops below the established baseline, the system flags the change as a feeding irregularity that requires inspection. In this way, the system helps improve feeding consistency, reduce feed waste, and provide caretakers with a clearer basis for identifying pens that need attention.
\\

\noindent\textit{Animal Health and Biosecurity}

The proposed system provides caretakers with continuous records of feed consumption, water level, and air-quality conditions in each monitored pen. These records allow caretakers to compare current pen conditions with previous readings instead of relying only on manual observation. When the system detects abnormal feed consumption, low water level, or air-quality readings beyond the set threshold, it generates an alert for caretaker inspection.

Through automated alerts and structured data records, caretakers can respond more promptly to feeding irregularities, water supply issues, or unsafe air-quality conditions. These alerts provide caretakers with timestamped information on the type and location of the detected abnormal condition, allowing them to inspect the affected pen based on recorded sensor readings rather than manual observation alone. In this way, the system supports organized record-keeping and consistent monitoring of feeding, watering, and environmental conditions.
\\

\noindent\textit{The Philippine Swine Industry}

With approximately 71.1\% of the Philippine hog inventory managed by smallholder backyard farmers, the local swine sector remains highly vulnerable to disruptions caused by disease and inefficient farm management. A reliable, centralized automated system could help stabilize domestic pork production, reduce dependence on imported meat, which reached a record 1.64 million metric tons in 2025, and support the national goal of rebuilding the country's swine population to pre-ASF levels.
\\

\noindent\textit{Field of Agricultural Automation and Embedded Systems}

This study contributes to the growing body of research on IoT and precision agriculture applied to livestock management. By demonstrating a fully hardwired, sensor-integrated feeding and monitoring system designed specifically for the Philippine piggery context, the project provides a practical reference for future engineering work in agricultural automation, particularly in environments where wireless reliability may be a concern.
\\

\noindent\textit{Future Researchers}

The system's data logging and analytics capabilities generate structured, timestamped records of feed consumption, water intake, and environmental conditions. This dataset can serve as a foundation for future studies on swine behavioral patterns, machine learning-based health prediction models, and the broader integration of precision livestock farming technologies in developing agricultural economies.
